% ============================================================================================================
\chapter{Limites restantes et travaux à mener}\label{ch:limites}
% ============================================================================================================

\section{Ce que la v0.7 ne fait pas encore}

\begin{enumerate}
\item \textbf{Les effets distributifs restent statiques.} Les gagnants et les perdants sont calculés sur les revenus
de l'année de la microsimulation, avant les effets macroéconomiques (emploi, salaires, prix). Une mesure qui crée des
emplois améliore la situation de ménages qui n'apparaissent pas comme gagnants dans ce calcul. L'inverse vaut pour une
mesure qui détruit des emplois. C'est la limite la plus importante pour l'acceptabilité.
\item \textbf{Les comportements des entreprises sont agrégés.} Le modèle ne distingue pas les entreprises par taille ;
les branches sont distinguées pour la production, l'emploi et l'investissement, mais les réactions à l'IS, à la
trésorerie et à la TVA sont uniformes. Le transfert de bénéfices et l'additionnalité de la R\&D sont, en réalité,
concentrés sur les grands groupes ou les PME.
\item \textbf{Pas de politique monétaire propre, taux souverains exogènes.} La BCE réagit à la zone euro, pas à la
France ; le module M11 représente une partie de cet effet. Une dégradation de la dette ne relève pas les taux
français dans le scénario central (pas de prime de risque endogène).
\item \textbf{Pas de bilans patrimoniaux.} Pas d'effets de richesse, de désendettement privé, de prix immobiliers.
\item \textbf{Les leviers d'offre sont fragiles.} Les effets de long terme de la R\&D, de la formation, de la durée du
travail et de l'âge de départ passent par le PIB potentiel ; ils dépendent d'élasticités dont la transposition à
grande échelle est incertaine.
\item \textbf{L'anticipation est une hypothèse.} Le paramètre $\theta$ n'est pas estimé ; il est choisi pour que le
modèle se comporte de façon cohérente après un choc d'offre. Les résultats de long terme des mesures qui modifient le
potentiel en dépendent.
\item \textbf{L'efficacité de la dépense publique n'est pas représentée.} Dans le modèle, une économie de dépense
réduit la demande de la même façon, qu'elle supprime un service ou un gaspillage. Les gains d'efficacité ne peuvent
donc pas être distingués des coupes.
\item \textbf{La référence est conventionnelle.} La trajectoire « politique inchangée » conduit à une dette de près de
180 \% du PIB en 2037 et de 300 \% en 2047, avec des taux de marché à 4,9 \%. Elle sert d'étalon pour mesurer les
écarts, non de prévision.
\item \textbf{L'optimisation repose sur un métamodèle linéaire.} Les interactions entre mesures sont vérifiées a
posteriori par le modèle complet pour les stratégies retenues, pas pour tout le front de Pareto.
\end{enumerate}

\section{Travaux prioritaires pour les phases suivantes}

\begin{enumerate}
\item Microsimulation dynamique : faire évoluer les revenus des ménages avec l'emploi, les salaires et les prix du
modèle macroéconomique, pour publier des gagnants et perdants « après effets de second tour ».
\item Démographie des entreprises : distinguer TPE, PME, ETI et grands groupes (données ouvertes Ésane, FARE,
Insee).
\item Prime de risque souveraine endogène et politique monétaire de la zone euro.
\item Estimation, et non calibration, des paramètres de la fermeture et de l'anticipation (méthode des moments sur
les chocs historiques).
\item Bilans sectoriels (comptes de patrimoine de l'Insee et de la Banque de France).
\item Représentation de l'efficacité de la dépense publique (indicateurs de production des services publics).
\item Énergie : un levier de politique (prix réglementés, contrats de long terme, subventions) plutôt qu'un prix
exogène.
\end{enumerate}
